\documentclass[12pt,a4paper]{article}
\usepackage[utf8]{inputenc}
\usepackage[T1]{fontenc}
\usepackage[margin=2.5cm,headheight=16pt,headsep=18pt]{geometry}
\usepackage{graphicx}
\usepackage{fancyhdr}
\usepackage{booktabs}
\usepackage{tabularx}
\usepackage{amsmath}
\usepackage{tikz}
\usepackage{listings}
\usepackage{inconsolata}
\usepackage{caption}
\usepackage[section]{placeins}
\usepackage{needspace}
\usepackage{float}
\usepackage[normalem]{ulem}
\usepackage[hidelinks]{hyperref}
\graphicspath{{docs/}{./}}

\renewcommand{\contentsname}{Daftar Isi}
\renewcommand{\appendixname}{Lampiran}
\renewcommand{\tablename}{Tabel}
\renewcommand{\figurename}{Gambar}
\setcounter{tocdepth}{3}
\setlength{\parindent}{1.5em}
\setlength{\parskip}{0pt}
\lstset{language=Python,basicstyle=\ttfamily\footnotesize,keywordstyle=\bfseries,emphstyle=\bfseries,breaklines=true,frame=single,showstringspaces=false,columns=fullflexible,aboveskip=3pt,belowskip=3pt,floatplacement=htbp,float=htbp,float=true}
\makeatletter
\setlength{\@fpsep}{10pt}
\makeatother
\setlength{\textfloatsep}{10pt}
\setlength{\intextsep}{10pt}
\captionsetup{font=small,labelfont=bf,labelsep=period,skip=6pt}
\raggedbottom
\clubpenalty=10000
\widowpenalty=10000
\displaywidowpenalty=10000
\renewcommand{\topfraction}{0.9}
\renewcommand{\bottomfraction}{0.8}
\renewcommand{\textfraction}{0.08}
\renewcommand{\floatpagefraction}{0.75}
\AddToHook{cmd/subsection/before}{\FloatBarrier\Needspace{7\baselineskip}}
\AddToHook{cmd/subsubsection/before}{\FloatBarrier\Needspace{5\baselineskip}}
\renewcommand{\lstlistingname}{Cuplikan}

\pagestyle{fancy}
\fancyhf{}
\fancyhead[L]{IF3170 -- Inteligensi Artifisial}
\fancyhead[R]{\thepage}
\renewcommand{\headrulewidth}{0.4pt}
\renewcommand{\footrulewidth}{0pt}

\begin{document}

\hypersetup{pageanchor=false}
\begin{titlepage}
    \centering
    {\Huge\bfseries Laporan Tugas Besar 2\par}
    \vspace{0.5cm}
    {\LARGE\scshape IF3170 Inteligensi Artifisial\par}
    \vspace{0.3cm}
    {\Large\scshape Klasifikasi Pembatalan Pemesanan Penginapan\par}
    \vspace{0.2cm}
    {\Large\scshape dengan Algoritma Supervised Learning\par}
    \vspace{0.3cm}
    {\Large\scshape Tahun Akademik 2026--2027\par}
    \vfill
    \includegraphics[width=7.5cm,keepaspectratio]{logo-itb.png}\par
    \vfill
    {\large\bfseries Disusun oleh:\par}
    \vspace{0.3cm}
    {\large LalapKaoBole\par}
    \vspace{0.3cm}
    {\large Nicholas Wise Saragih Sumbayak -- 13524037\par}
    {\large Nathan E C Marpaung -- 13524062\par}
    {\large Kurt Mikhael Purba -- 13524065\par}
    {\large Aurelia Jennifer Gunawan -- 13524089\par}
    {\large Amanda Aurellia Salsabila -- 13524131\par}
    \vfill
    {\large\scshape Laboratorium Inteligensi Artifisial\par}
    {\large\scshape Program Studi Teknik Informatika\par}
    {\large\scshape Sekolah Teknik Elektro dan Informatika\par}
    {\large\scshape Institut Teknologi Bandung\par}
\end{titlepage}

\pagenumbering{arabic}
\hypersetup{pageanchor=true}
{\small\linespread{0.92}\selectfont\tableofcontents}
\clearpage

\section{Pendahuluan}
\subsection{Latar Belakang}

\subsection{Rumusan Persoalan dan Tujuan}

\subsection{Deskripsi Dataset}
\subsubsection{Sumber dan Struktur Dataset}
\subsubsection{Atribut dan Tipe Data}
\subsubsection{Variabel Target}
\subsubsection{Eksplorasi Data}

\subsection{Batasan Persoalan}

\section{Dasar Teori}
\subsection{Supervised Learning}
\subsection{Klasifikasi}
\subsection{K-Nearest Neighbors}
\subsection{Decision Tree Learning}
\subsubsection{Iterative Dichotomiser 3 (ID3)}
\subsubsection{C4.5}
\subsubsection{Classification and Regression Tree (CART)}
\subsection{Logistic Regression}
\subsubsection{Regularisasi L1}
\subsubsection{Regularisasi L2}
\subsubsection{Elastic Net}
\subsection{Support Vector Machine}
\subsubsection{Kernel Linear}
\subsubsection{Kernel Polinomial}
\subsubsection{Kernel Gaussian}
\subsection{Bayesian Network (Naive Bayes)}
\subsection{Validasi dan Metrik Evaluasi}

\section{Pembahasan dan Implementasi}
\subsection{Arsitektur Aplikasi}
\subsubsection{Struktur Aplikasi}
\subsubsection{Tech Stack}
\subsection{Deskripsi Fungsi dan Kelas}
\subsubsection{K-Nearest Neighbors}
\subsubsection{Decision Tree Learning}
\subsubsection{Logistic Regression}
\subsubsection{Support Vector Machine}
\subsubsection{Bayesian Network (Naive Bayes)}
\subsection{Data Cleaning}
\subsubsection{Penanganan Nilai Hilang}
\subsubsection{Penanganan Data Duplikat}
\subsubsection{Penanganan Data Tidak Valid}
\subsection{Data Transformation}
\subsubsection{Encoding Fitur Kategorikal}
\subsubsection{Normalisasi dan Standardisasi Fitur Numerik}
\subsubsection{Penanganan Ketidakseimbangan Kelas}
\subsection{Feature Selection}
\subsection{Dimensionality Reduction}
\subsection{Implementasi Pembanding dengan scikit-learn}
\subsection{Penyimpanan dan Pemuatan Model}

\section{Eksperimen dan Analisis}
\subsection{Rancangan Eksperimen}
\subsubsection{Data Training dan Data Test}
\subsubsection{Pembagian Data dan Validasi}
\subsubsection{Konfigurasi Parameter}
\subsubsection{Metrik Evaluasi}
\subsection{Hasil Eksperimen dan Visualisasi}
\subsubsection{K-Nearest Neighbors}
\subsubsection{Decision Tree Learning}
\subsubsection{Logistic Regression}
\subsubsection{Support Vector Machine}
\subsubsection{Bayesian Network (Naive Bayes)}
\subsection{Perbandingan Implementasi from-scratch dengan scikit-learn}
\subsubsection{Perbandingan Metrik}
\subsubsection{Analisis Waktu Eksekusi}
\subsubsection{Analisis Visualisasi}
\subsection{Kaggle Submission}

\clearpage
\section{Kesimpulan dan Saran}
\subsection{Kesimpulan}
\subsection{Saran}

\pagebreak
\section{Pembagian Tugas Tiap Anggota Kelompok}
\begin{table}[ht]
    \centering
    \renewcommand{\arraystretch}{1.5}
    \begin{tabularx}{\textwidth}{@{}p{6.1cm}lX@{}}
        \toprule
        \textbf{Nama Lengkap} & \textbf{NIM} & \textbf{Tugas/Kontribusi} \\
        \midrule
        Nicholas Wise Saragih Sumbayak & 13524037 & \\
        Nathan E C Marpaung & 13524062 & \\
        Kurt Mikhael Purba & 13524065 & \\
        Aurelia Jennifer Gunawan & 13524089 & \\
        Amanda Aurellia Salsabila & 13524131 & \\
        \bottomrule
    \end{tabularx}
    \caption{Pembagian tugas anggota kelompok}
\end{table}

\clearpage
\section{Referensi}

\clearpage
\section{Lampiran}
\subsection{Form Penggunaan \textit{Generative} AI}
\begin{center}
    {\large\bfseries SURAT PERNYATAAN PENGGUNAAN AI\par}
\end{center}
\vspace{-0.3em}
\noindent Dengan ini, \sout{saya}/kami*:

\begin{tabular}{@{}l@{\hspace{1.2em}}l@{\hspace{0.8em}}p{0.62\textwidth}@{}}
    Nama/NIM & : & Nicholas Wise Saragih Sumbayak/13524037, Nathan E C Marpaung/13524062, Kurt Mikhael Purba/13524065, Aurelia Jennifer Gunawan/13524089, Amanda Aurellia Salsabila/13524131 \\
    Fakultas/Sekolah & : & STEI-K \\
    Kode Mata Kuliah & : & IF3170 \\
    Nama Mata Kuliah & : & Inteligensi Artifisial \\
    Judul Tugas/Proposal & : & Tugas Besar 2 \\
\end{tabular}

\vspace{0.6em}
\noindent Menyatakan bahwa \sout{saya}/kami* \underline{menggunakan/tidak menggunakan}* AI dalam pengerjaan maupun penyusunan luaran tugas pada mata kuliah yang tertulis di atas.

\vspace{0.4em}
\noindent Jika Ya, maka alat AI/kecerdasan buatan yang digunakan adalah:

\vspace{0.4em}
\begin{table}[H]
    \centering
    \small
    \renewcommand{\arraystretch}{1.15}
    \begin{tabularx}{\textwidth}{|>{\raggedright\arraybackslash}X|c|>{\centering\arraybackslash}p{2.6cm}|}
        \hline
        \multicolumn{1}{|c|}{\textbf{Lingkup Pekerjaan}} & \textbf{Digunakan?} & \textbf{Tingkat Penggunaan} \\
        \hline
        \textbf{Pemeriksaan Ejaan:} menggunakan tools seperti Grammarly, Paperpal, DeepL, Quillbot, Grammarbot, LanguageTool, ProWritingAid, ChatGPT, Google Gemini, atau sejenisnya. & & \\
        \hline
        \textbf{Pembuatan Teks:} menggunakan tools seperti ChatGPT, Gemini, Paperpal, Copilot, GrammarlyGO, WordAI, WriteSonic, Jasper, Jenni AI, atau sejenisnya. & & \\
        \hline
        \textbf{Bantuan Diskusi dalam Penyusunan Konten:} menggunakan tools seperti ChatGPT, Gemini, Copilot, Perplexity, Jenni AI, Zoom-Companion, atau sejenisnya. & & \\
        \hline
        \textbf{Pembuatan Gambar, Video, dan/atau Grafik:} menggunakan tools seperti Craiyon, DALL-E, Midjourney, Stable Diffusion, Microsoft Designer, Gemini, Canva AI, atau sejenisnya. & & \\
        \hline
        \textbf{Lainnya, Jelaskan:}\vspace{1.6em} & & \\
        \hline
    \end{tabularx}
\end{table}

\noindent \sout{Saya}/Kami* juga menyatakan bahwa setiap penggunaan AI/kecerdasan buatan yang dilakukan pada mata kuliah tersebut telah dinyatakan di dalam surat ini.

\vspace{0.5em}
\hfill\begin{minipage}{6.5cm}
    \raggedleft
    Bandung, \rule{3cm}{0.4pt}\par
    \vspace{1.8cm}
    \rule{5cm}{0.4pt}\par
    Nama dan NIM
\end{minipage}

\end{document}
